\documentclass[11pt]{article}
\usepackage[margin=2.6cm]{geometry}
\usepackage{amsmath,amssymb,amsthm}
\usepackage{enumitem}
\setlist{itemsep=1pt,topsep=3pt}

\newcommand{\C}{\mathcal{C}}
\newcommand{\Path}{\mathrm{Path}}
\newcommand{\id}{\mathrm{id}}
\newcommand{\D}{\mathsf{D}}
\newcommand{\Lam}{\Lambda}

\theoremstyle{plain}
\newtheorem{theorem}{Theorem}[section]
\newtheorem{lemma}[theorem]{Lemma}
\newtheorem{proposition}[theorem]{Proposition}
\newtheorem{corollary}[theorem]{Corollary}
\theoremstyle{definition}
\newtheorem{definition}[theorem]{Definition}
\theoremstyle{remark}
\newtheorem{remark}[theorem]{Remark}

\title{$\D A$ as a limit of the $\D_NA$\\[3pt]
\large strict limit, homotopy limit, and a tower without fibrations}
\author{}
\date{22 October 2026}

\begin{document}
\maketitle

\noindent\textbf{Mathematical attribution. All mathematical results
concerning the $\D$ modality and the $\D_n$ modalities presented here are
due to Christian Sattler.} This includes the constructions, proofs,
finite approximations, and all results about their tower and its limits.
This note records their exposition and formalisation in Unit.

\begin{abstract}
\noindent
For a type $A$ of the poset model, $\D_NA$ is the localisation $\D A$
with the families indexed only by the simplices of dimension at most $N$.
We present three results of Christian Sattler. $\D A$ is the strict limit
of the tower $(\D_NA)_N$. If $A$ is fibrant, $\D A$ is also its homotopy
limit: the comparison map from $\D A$ to the internal limit, a type built
with $\Pi$, $\Sigma$ and $\Path$, is a homotopy equivalence. Finally, the
maps $\D_{N+1}A\to\D_NA$ need not be fibrations, even when $A$ is fibrant,
so the second fact is not a consequence of the first. All three are checked
in \texttt{core}.
\end{abstract}

%======================================================================
\section{Setting}\label{sec:tower}

\paragraph{The poset model.} $\C$ is the category of finite non-empty
posets and monotone maps, and $\Delta\subseteq\C$ the full subcategory of
the linear posets $[n]=\{0<\dots<n\}$. A map of $\Delta$ is \emph{strict}
if it is strictly monotone, and \emph{surjective} if it is onto. Every map
$g$ of $\Delta$ factors uniquely as $g=c\sigma$ with $c$ strict and
$\sigma$ surjective (Eilenberg--Zilber), and the same holds for a map
$g:[k]\to P$ into a poset, with $c:[r]\to P$ injective. A context $\Gamma$
is a presheaf on $\C$, and a type $A$ over $\Gamma$ a presheaf on its
category of elements; we write $A(X,\rho)$ for the fibre at $\rho\in\Gamma(X)$
and $a\,g$ for restriction. Every poset is identified with its nerve, so
that $[m]\times[1]^k$ is $\Delta[m]\times\Delta[1]^k$. We write $i^*A$ for
the restriction of $A$ to $\Delta$, a simplicial set over $i^*\Gamma$, and
$\partial[m]$ for the union of the faces of $[m]$.

A \emph{path} in $A$ at $(X,\rho)$ from $a$ to $b$ is an element of $A$
at $X\times[1]$ over $\rho\pi_1$ whose restrictions to $X\times\{0\}$ and
$X\times\{1\}$ are $a$ and $b$. A family of paths natural in $X$ is a
section of the type $\Path_A$. A \emph{horn filling operation} on a
simplicial set $T$ chooses fillers of horns $\Lam^k[n]\to T$ such that the
filler of the restriction of a horn along a face map is the restriction of
the filler, whenever both are horns. We use two facts about fibrant types
of the poset model.
\begin{enumerate}
\item[(F1)] If $A$ is fibrant, $i^*A$ has a horn filling operation.
\item[(F2)] If $i^*B$ has a horn filling operation, the type $\D B$ defined
  below is fibrant, naturally in $\Gamma$.
\end{enumerate}

\begin{definition}
Let $A$ be a type over $\Gamma$. An element of $\D A$ at $(X,\rho)$ is a
family
\[
u(f)\in A([k],\rho f),\qquad f:[k]\to X,
\]
with $u(f)\,g=u(fg)$ for every strict $g:[j]\to[k]$. Restriction along
$X'\to X$ is precomposition. For $N\ge-1$, an element of $\D_NA$ at
$(X,\rho)$ is the same with only the maps $f:[k]\to X$ with $k\le N$. The
map $r:\D_{N+1}A\to\D_NA$ forgets the values in dimension $N+1$, and
$v_N:\D A\to\D_NA$ forgets the values above $N$.
\end{definition}

The important point is that the elements of $\D A$ and of $\D_NA$ are
natural along the strict maps only. Their values on degenerate simplices
are not determined by the values on the non-degenerate ones.

\paragraph{The coskeleton.} An element of $K_NA$ at $(X,\rho)$ is a
family $c(f)\in A([k],\rho f)$, for $k\le N$ and $f:[k]\to X$, with
$c(f)\,g=c(fg)$ for every $g:[j]\to[k]$ with $j\le N$. So $i^*K_NA$ is the
$N$-coskeleton of $i^*A$, and $\eta(a)(f)=af$ defines $\eta:A\to K_NA$.

\begin{lemma}\label{lem:cosk}
\leavevmode
\begin{enumerate}
\item For $m\le N$, $\eta:A([m],\tau)\to K_NA([m],\tau)$ is a bijection,
  with inverse $c\mapsto c(\id)$.
\item If $i^*A$ has a horn filling operation, so has $i^*K_NA$.
\end{enumerate}
\end{lemma}

\begin{proof}
(1) The identity of $[m]$ is one of the indices, and $c(g)=c(\id)\,g$.

(2) Let $h$ be a horn of $i^*K_NA$ on $\Lam^k[m+1]$. If $m\le N$, its faces
are, by (1), elements of $A([m])$ forming a horn of $i^*A$; its filler $a$
gives the filler $\eta(a)$. If $m>N$, a map $[j]\to[m+1]$ with $j\le N$
misses at least two vertices, hence a vertex other than $k$, so it lies in
the horn, and the filler is $h$ read at these maps. Both cases commute with
restriction to faces.
\end{proof}

\begin{proposition}\label{prop:DN}
$\D_NA=\D(K_NA)$. In particular $\D_NA$ is fibrant when $A$ is.
\end{proposition}

\begin{proof}
Given $u\in\D(K_NA)(X,\rho)$, put $w(f)=u(f)(\id)$ for $f:[k]\to X$, $k\le
N$; it is natural along strict maps. Conversely, given $w\in\D_NA(X,\rho)$,
put $u(f)(g)=w(fc)\,\sigma$, where $g=c\sigma$ is the factorisation of
$g:[j]\to[k]$ with $c$ strict and $\sigma$ surjective. Two such
factorisations are related by a strict map, which gives the naturality of
$u(f)$ along every map and of $u$ along strict maps. Below $N$ the value
$w(fc)\sigma$ is $w(f)\,c\sigma=w(f)g$, so the two maps are inverse. The
second statement follows from (F1), Lemma~\ref{lem:cosk}(2) and (F2).
\end{proof}

%======================================================================
\section{The strict limit}\label{sec:strict}

\begin{proposition}\label{prop:strict}
$\D A$ is the limit of the tower
$\cdots\to\D_{N+1}A\xrightarrow{r}\D_NA\to\cdots\to\D_{-1}A=1$, with the
cone $(v_N)_N$.
\end{proposition}

\begin{proof}
A compatible sequence $(u_N)_N$, $u_N\in\D_NA(X,\rho)$, defines the family
$u(f)=u_k(f)$ for $f:[k]\to X$. Compatibility gives $u_k(f)=u_N(f)$ for
$k\le N$, so $u$ is natural along every strict map, since $u_N$ is for
$N$ large. This $u$ is the unique element of $\D A(X,\rho)$ with
$v_N(u)=u_N$ for all $N$, and the construction commutes with restriction.
\end{proof}

%======================================================================
\section{The maps of the tower are not fibrations}\label{sec:notfib}

\begin{proposition}\label{prop:notfib}
There is a fibrant type $A$ such that $r:\D_1A\to\D_0A$ has no fibration
structure.
\end{proposition}

\begin{proof}
Let $A$ be the constant type with value $\{0,1\}$ over the empty context,
which is fibrant. Then $A([k])=\{0,1\}$ for every $k$, and every
restriction is the identity. An element of $\D_0A$ at $X$ is an arbitrary
function from the points of $X$ to $\{0,1\}$, since the only strict map
$[0]\to[0]$ is the identity. An element of $\D_1A$ at $X$ is such a
function together with a value on each map $[1]\to X$. Its value on an
edge $e$ restricts, along the two strict maps $[0]\to[1]$, to its values on
the two ends of $e$. So the ends of every edge get the same value.

Take $X=[0]$. The element $w=(0)$ of $\D_1A([0])$ lies over $a=0$. The
element $\ell$ of $\D_0A([1])$ with $\ell(0)=0$ and $\ell(1)=1$ is a line
of $\D_0A$ from $r(w)$ to $1$. A lift of $\ell$ starting at $w$ is an
element of $\D_1A([1])$ over $\ell$. It must give a value to the identity
of $[1]$, an edge from $0$ to $1$, and there is none. So $r$ does not even
have the path lifting property. A fibration structure gives fillers for
all open boxes over $[0]\times[1]$, and in particular such a lift.
\end{proof}

\begin{remark}
The same argument applies to $r:\D_{N+1}A\to\D_NA$ for every $N\ge0$, for a
fibrant $A$ whose restriction to $\Delta$ has an $N$-sphere with no
filler. Lifting a line of $\D_NA$ over $[N]$ asks for $(N+1)$-simplices of
$i^*A$ with a prescribed boundary. This is the coskeletal tower of a Kan
complex: it converges strictly although its maps are not Kan fibrations.
\end{remark}

So the strict limit of Proposition~\ref{prop:strict} is not the limit of a
tower of fibrations, and it need not compute the homotopy limit. This is
the content of the next section.

%======================================================================
\section{The homotopy limit}\label{sec:hlim}

The \emph{internal limit} of the tower is the type
\[
L=\Sigma\,(u:\textstyle\prod_N\D_NA)\,\prod_N\Path_{\D_NA}(r\,u_{N+1},u_N),
\]
where the products range over the constant type of natural numbers. The
\emph{comparison} $\varphi:\D A\to L$ sends $u$ to $((v_Nu)_N,\text{the
constant paths})$.

If $A$ is fibrant, so is $L$: each $\D_NA$ is fibrant
(Proposition~\ref{prop:DN}), and fibrant types are closed under $\Path$,
$\Sigma$ and products over the constant type of natural numbers, the
latter componentwise.

\begin{theorem}\label{thm:hlim}
If $A$ is fibrant, $\varphi$ is a homotopy equivalence: there is a map
$\psi:L\to\D A$, natural in the context, with paths
$\varphi\psi\simeq\id_L$ and $\psi\varphi\simeq\id_{\D A}$ in the model.
\end{theorem}

The rest of this section proves the theorem. Write $C=i^*A$ and
$C_N=i^*K_NA$, the $N$-coskeleton of $C$. We identify a poset with its
nerve, so that $[m]\times[1]^k$ is $\Delta[m]\times\Delta[1]^k$.

\subsection{Tools}

\begin{lemma}[Boxes]\label{lem:box}
Let $T$ be a simplicial set with a horn filling operation, $P$ a finite
poset, $S\subseteq P$ a simplicial subset and $e\in\{0,1\}$. Every map
$P\times\{e\}\cup S\times\Delta[1]\to T$ extends to $P\times\Delta[1]$,
with a choice of extension compatible with restriction along injective
maps $P'\to P$.
\end{lemma}

\begin{proof}
\emph{Prisms over a simplex.} A problem at $n$ is a map
$b:[n]\times\{e\}\cup Q\times\Delta[1]\to T$ with $Q\subseteq[n]$; its
restriction along $d_i$ is a problem at $n-1$. By recursion on $n$ we
choose a filler $F(b)$ with $F(b)\,d_i=F(b\,d_i)$. If $Q$ contains the
identity, $F(b)=b$. Otherwise the fillers $F(b\,d_i)$ agree on common
faces, and $\partial[n]\times\Delta[1]\cup[n]\times\{e\}$ is given. The
non-degenerate $(n+1)$-simplices of $[n]\times[1]$ are the shuffles
$P_k=((0,0),\dots,(k,0),(k,1),\dots,(n,1))$, $0\le k\le n$. The face $d_k$
of $P_k$ is the face $d_k$ of $P_{k-1}$, the face $d_{k+1}$ of $P_k$ is the
face $d_{k+1}$ of $P_{k+1}$, and the other faces lie in
$\partial[n]\times\Delta[1]$, apart from $d_{n+1}P_n=[n]\times\{0\}$ and
$d_0P_0=[n]\times\{1\}$. For $e=0$ fill $P_n,\dots,P_0$ in this order, each
as the filler of a horn on $\Lam^k[n+1]$; for $e=1$ fill $P_0,\dots,P_n$,
each on $\Lam^{k+1}[n+1]$.

\emph{Prisms over a poset.} For an injective $c:[n]\to P$ let $b_c$ be the
problem read along $c$, with $Q=c^*S$. Then $b_c\,d_i=b_{cd_i}$, so
$F(b_c)\,\mu=F(b_{c\mu})$ for every injective $\mu$. A simplex $g:[W]\to P$
factors uniquely as $g=c\sigma$, and we put $F(g,j)=F(b_c)(\sigma,j)$. If
$k:[W']\to[W]$, the factorisation of $gk$ is $(c\mu)\sigma'$ with $\mu$
injective and $\mu\sigma'=\sigma k$, so $F(gk,jk)=F(g,j)\,k$. If $g\in S$
then $c\in S$, the identity is in $c^*S$, and $F(g,j)$ is the given value;
the same holds at $j=e$.
\end{proof}

We apply the lemma with $P=[m]\times Y$ and
$S=\partial[m]\times Y\cup[m]\times Q$, for a poset $Y$ and a simplicial
subset $Q\subseteq Y$. The \emph{lid} is the value at $e$, the \emph{sides}
are the values on $\partial[m]\times Y$, and the \emph{end data} are the
values on $[m]\times Q$. For the squares, $Y=[1]_t$ and $Q=\{0,1\}$. For the
cubes, $Y=[1]_t\times[1]_v$ and $Q$ is the boundary of the square: four
pieces at $t=0$, $t=1$, $v=0$, $v=1$, which must agree at the four corners.
The fillers are uniform along the face maps, so no excluded middle is used.

\paragraph{Readings.} Let $P$ be a finite poset and $\phi:P\to X$. An
element $z$ of $\D A$ (or of $\D_NA$, if the simplices of $P$ have
dimension at most $N$) at $X$ is \emph{read} along $\phi$ as the map
$P\to C$ given on an injective chain $c:[r]\to P$ by $z(\phi c)$, and
extended to all simplices by the Eilenberg--Zilber factorisation $g=c\sigma$,
with $c$ injective and $\sigma$ surjective. The reading is natural along
all maps of $\Delta$, although $z$ is natural along strict maps only.

Let $u\in\D A(X)$ and $f:[m]\to X$, and put $U_f=u(f)\in C_m$. Reading $u$
along $f\pi_1:[m]\times[k]\to X$ gives a map $R_f^{[k]}$. For $k=0$ it is
the simplex $U_f$. We write $R_f=R_f^{[1]}$, a prism with both ends $U_f$.

\begin{lemma}\label{lem:R}
$R_f$ need not be the constant prism at $U_f$. There is a square $Q_f$ over
$[m]\times[1]_s\times[1]_v$ with
\[
Q_f(s,0)=R_f(s),\qquad Q_f(1,v)=R_f(v),\qquad Q_f(0,v)=R_f(v),\qquad
Q_f(s,1)=U_f .
\]
Readings commute with the restriction to faces: $R_f$, $Q_f$ and $U_f$
restricted along $d_i\times\id$ are $R_{fd_i}$, $Q_{fd_i}$ and $U_{fd_i}$.
\end{lemma}

\begin{proof}
On the edge $(i,0)\to(i,1)$ of $[m]\times[1]$, $R_f$ is $u$ at a degenerate
simplex of $X$. The naturality of $u$ along strict maps says nothing about
this value, which is not determined by $u(f)$.

Let $\beta:[1]\times[1]\to[2]$ be the monotone map with $\beta(s,0)=s$,
$\beta(1,v)=1+v$, $\beta(0,v)=2v$ and $\beta(s,1)=2$. Put
$Q_f=R_f^{[2]}\circ(\id\times\beta)$. The edges $v=0$, $s=1$ and $s=0$ of
$\beta$ are the three strict maps $[1]\to[2]$. Reading along a strict
inclusion $P'\to P$ is reading on $P'$, since the injective chains of $P'$
are injective chains of $P$. So these three edges of $Q_f$ are $R_f$. The
edge $v=1$ is $R_f^{[2]}$ precomposed with $[m]\to[m]\times\{2\}$, which is
again a strict inclusion, so it is the reading of $u$ along $f$, which is
$U_f$. Note that we precompose an already extended map. We do not claim
that reading commutes with projections.

The last statement is the same argument for the strict inclusions
$d_i\times\id$, together with $u(f)\,d_i=u(fd_i)$.
\end{proof}

\subsection{The map $\psi$}

Let $\ell=(u_N,p_N)\in L(X)$, so $p_N$ is an element of $\D_NA$ at
$X\times[1]$ from $r\,u_{N+1}$ to $u_N$. For $f:[m]\to X$, let $\sigma_N(f)$
be the reading of $u_N$ along $f$ and $p_N(f)$ the reading of $p_N$ along
$f\times[1]$. We build, by recursion on $m$ and, for each $f$, on the
stage $N\ge m+1$:
\begin{itemize}
\item an $m$-simplex $\psi(f)$ of $C$;
\item prisms $H^N(f)$ over $[m]\times[1]_t$ from $\psi(f)$ ($t=0$) to
  $\sigma_N(f)$ ($t=1$);
\item squares $\Theta_N(f)$ over $[m]\times[1]_t\times[1]_s$ with
  $\Theta_N(t,1)=H^N(t)$, $\Theta_N(t,0)=H^{N+1}(t)$,
  $\Theta_N(0,s)=\psi(f)$ and $\Theta_N(1,s)=p_N(f)(s)$.
\end{itemize}
Each of these restricts along $d_i$ to the same datum at $fd_i$.
\begin{enumerate}
\item $H^{m+1}(f)$ fills the box over $[m]\times[1]_t$ with lid
  $\sigma_{m+1}(f)$ at $t=1$ and sides $H^{m+1}(fd_i)$. Its end at $t=0$ is
  $\psi(f)$.
\item $\Theta_N(f)$ fills the box over $[m]\times[1]_t$, line $s$, with lid
  $H^N(f)$ at $s=1$, end data $\psi(f)$ at $t=0$ and $p_N(f)$ at $t=1$, and
  sides $\Theta_N(fd_i)$. Its end at $s=0$ is $H^{N+1}(f)$.
\end{enumerate}
The face $fd_i$ has dimension $m-1$, so its data exist for $N\ge m$, and
in particular at the stages $N\ge m+1$ used at $f$. All these boxes lie in
dimension at most $m+2\le N+1$. They live in $C$: the stage only says which
of $u_N$ and $p_N$ is read. Then $\psi(\ell)=(f\mapsto\psi(f))$ is natural
along the face maps, hence along all strict maps, and is an element of
$\D A(X)$.

The value of $\psi(\ell)$ in dimension $m$ is taken at stage $m+1$, never
lower. In particular $u_0$ is never used, so the obstruction of
Section~\ref{sec:notfib} does not arise. The squares $\Theta_N$ are part of
the data and are built by box filling, not by concatenating the paths
$p_N$: a concatenation would not commute strictly with the faces.

\subsection{The path $\varphi\psi\simeq\id_L$}

A path in $L$ from $\varphi\psi(\ell)$ to $\ell$ consists of lines $l_N$ of
$\D_NA$ from $v_N\psi(\ell)$ to $u_N$, and squares relating $r\,l_{N+1}$,
$l_N$, the constant line and $p_N$. The line $l_N$ is read off $H^N(f)$ at
the maps $(f,h):[k]\to X\times[1]$, $k\le N$, through $\langle\id,h\rangle$:
its value at $(f,h)$ is $H^N(f)$ at $(\id,h)$, natural along strict maps
because $H^N$ is face compatible. For $k<N$ the prism
$H^N(f)$ exists. For $k=N$ it is built in $C_N$: $H^N(f)$ and
$\Theta_N(f)$ fill, in $C_N$, the box formed by $H^{N+1}(f)$, $p_N(f)$, the
constant line, and the sides $\Theta_N(fd_i)$. The squares are read off
$\Theta_N$ in the same way, and their edges are the required ones.

\subsection{The path $\psi\varphi\simeq\id_{\D A}$}

Let $u\in\D A(X)$ and run the construction on $\varphi(u)$. Its path
$p_N$ is constant, so its reading $p_N(f)$ is $R_f$ (Lemma~\ref{lem:R}).
Thus $H^N(f)$ goes from $\psi\varphi(u)(f)$ to $U_f$, and
\[
\Theta_N(t,1)=H^N(t),\quad\Theta_N(t,0)=H^{N+1}(t),\quad
\Theta_N(0,s)=\psi\varphi(u)(f),\quad\Theta_N(1,s)=R_f(s).
\]
We build, by recursion on $m$ and on $N\ge m+1$, a prism $G(f)$ over
$[m]\times[1]_t$ from $\psi\varphi(u)(f)$ to $U_f$ and squares
$\Xi^N(f)$ over $[m]\times[1]_t\times[1]_v$ with
\[
\Xi^N(t,0)=H^N(t),\quad\Xi^N(t,1)=G(t),\quad
\Xi^N(0,v)=\psi\varphi(u)(f),\quad\Xi^N(1,v)=R_f(v),
\]
all restricting along $d_i$ to the data at $fd_i$. We also keep the cubes
filled on the way.
\begin{enumerate}
\item \emph{Initial square.} $\Xi^{m+1}(f)$ fills the box over
  $[m]\times[1]_t$, line $v$, with lid $H^{m+1}(f)$ at $v=0$, end data
  $\psi\varphi(u)(f)$ at $t=0$ and $R_f(v)$ at $t=1$, and sides
  $\Xi^{m+1}(fd_i)$. Its end at $v=1$ is $G(f)$, which goes from
  $\psi\varphi(u)(f)$ to $R_f(1)=U_f$, with faces $G(fd_i)$.
\item \emph{Successor cube.} $\Xi^{N+1}(f)$ is the end at $s=0$ of the
  filler of the box over $[m]\times[1]_t\times[1]_v$, line $s$, with lid
  $\Xi^N(f)$ at $s=1$, sides the cubes at $fd_i$, and end data
  \[
  \begin{array}{ll}
  t=0: & \psi\varphi(u)(f),\\
  t=1: & Q_f(s,v),\\
  v=0: & \Theta_N(t,s),\\
  v=1: & G(t),\ \text{independent of }s.
  \end{array}
  \]
\end{enumerate}
The pieces agree at the corners: at $t=1,v=0$ both are $R_f(s)$; at
$t=1,v=1$ both are $U_f$; at $t=0$ all are $\psi\varphi(u)(f)$. At $s=1$
they are the edges of $\Xi^N(f)$: $Q_f(1,v)=R_f(v)$, $\Theta_N(t,1)=H^N(t)$.
At $s=0$ they give the required edges of $\Xi^{N+1}(f)$:
$Q_f(0,v)=R_f(v)$ and $\Theta_N(t,0)=H^{N+1}(t)$. On a face the end data
are those of the face, by Lemma~\ref{lem:R} and the face relations of
$\psi$, $H$, $\Theta$ and $G$. The sides agree with the lid because the
cube at $fd_i$ starts at $\Xi^N(fd_i)$, which is $\Xi^N(f)$ on that face.
The sides agree with each other where two faces meet, by the face
relations one dimension lower. All boxes are filled in $C$ itself: $R_f$
and $Q_f$ come from the full element $u$, so no truncation bound appears.

The prisms $G(f)$ are face compatible, so the values $G(f)(\id,h)$ for
$(f,h):[n]\to X\times[1]$ form an element of $\D A(X\times[1])$, natural
along the strict maps, from $\psi\varphi(u)$ to $u$.

\subsection{Naturality}

The data at $f$ are computed from the readings of the input along $f$ and
from the data at the faces of $f$. For $h:Y\to X$, the problem at $f$ for
the restriction of the input along $h$ is the problem at $hf$ for the
input. The fillers are chosen uniformly, so equal problems have equal
fillers, by induction on the dimension and on the stage. For
$\psi\varphi\simeq\id$ one needs in addition that the readings $R_f$ and
$Q_f$ of the restriction $h^*u$ are those of $u$ at $hf$, and that the
restriction of $\varphi(u)$ along $h$ is $\varphi(h^*u)$. So $\psi$ and
both homotopies are natural in $X$. A natural family of such cylinders is
a section of the path type, by the definition of paths in
Section~\ref{sec:tower}.
This proves Theorem~\ref{thm:hlim}.

\begin{remark}
No choice principle is used. The connectedness of the maps of the tower is
witnessed by box fillers of $C$ and of the $C_N$, chosen uniformly along
the face maps. The internal argument from towers of connected maps would
need dependent choice.
\end{remark}

%======================================================================
\appendix
\section{Correspondence with the formalisation}

\noindent Layer L9 of \texttt{core}, over the poset instance.

\medskip
\begin{tabular}{p{4.3cm}p{9.6cm}}
statement & files\\\hline
$\D_NA=\D(K_NA)$, fibrancy & CCosk, CCoskHorn, CLocDN, CLocDNC\\
Proposition~\ref{prop:strict} & CLocDN (Lim)\\
$L$, $\varphi$ & CPostLimit (Product, Tower, Cone, Post); $L$ fibrant: CPostFib\\
boxes & CBoxFill; over $[m]\times[1]$: CPostSq; over $[m]\times Y$: CPostBox; the
boundary of the square and its four pieces: CPostCube (Geo, E4, E4Cmp)\\
readings, Lemma~\ref{lem:R} & CPostMap (DR), CPostK; for $u\in\D A$:
CPostRead ($\beta$, Rd, Sq: $R$, $R2$, $Q$ and its edges, Face)\\
$\psi$ & CPostRec, CPostStep0/1/3, CPostInput\\
$\varphi\psi\simeq\id$ & CPostStep2, CPostLine, CPostSquare, CPostAssemble;
naturality CPostExt, CPostSqNat, CPostNat, CPostHomNat; the path:
CPostConvergence (RightInverse.homotopy)\\
$\psi\varphi\simeq\id$ & CPostInvIn (the recursion on $\varphi(u)$), CPostInv
($G$, $\Xi$, the cubes), CPostInvAsm (the cylinder, its ends), CPostInvNat
(naturality; LeftInverse.homotopy)
\end{tabular}

\medskip
Proposition~\ref{prop:notfib} is not formalised. The fibrancy of $L$ is
CPostFib (ProductFib: countable products; HTFib: the path types of the
edges; TowerFib, PostFib: $L$).

\end{document}
